\medskip

Soit la fonction $f$ définie sur $\R$ par 
$f(x) = \ln \left(1 + \text{e}^{-x}\right) + \dfrac14 x.$

On note $\mathcal{C}_f$ la courbe représentative de la fonction $f$ dans un repère orthonormé \Oij{} du plan.

\medskip

\textbf{Partie A}

\medskip

\begin{enumerate}
\item 
\begin{list}{\textbullet}{}
\item $\ds\lim _{x\to +\infty} \e^{-x}=0$ donc $\ds\lim _{x\to +\infty} \ln\left (1+\e^{-x}\right ) = \ln(1)=0$
\item $\ds\lim _{x\to +\infty} \dfrac{1}{4}x=+\infty$
\end{list}

Donc $\ds\lim _{x\to +\infty} f(x) = +\infty$.

\item On admet que la fonction $f$ est dérivable sur $\R$ et on note $f'$ sa fonction dérivée.
	\begin{enumerate}
		\item 
$\begin{gathered}f'(x)=\dfrac{-\e^{-x}}{1+\e^{-x}}+\dfrac{1}{4}
= \dfrac{\e^{x}\times \left (-\e^{-x}\right )}{\e^{x}\times \left (1+\e^{-x}\right )}+\dfrac{1}{4}
= \dfrac{-1}{\e^{x}+1} + \dfrac{\e^{x}+1}{4\left (\e^{x}+1\right )}
= \dfrac{-4+\e^{x}+1}{4\left (\e^{x}+1\right )}\\
\phantom{f'(x)}
= \dfrac{\e^{x}-3}{4\left (\e^{x}+1\right )}\end{gathered}$
				
		\item 
Sur $\R$, $f'(x)$ est du signe de $\e^{x}-3$.

$\e^{x}-3\geqslant 0
\iff
\e^{x}\geqslant 3
\iff
x \geqslant \ln(3)$

\begin{list}{\textbullet}{Donc:}
\item sur $\left ] -\infty\;;\;\ln(3)\strut\right [$, $f'(x)<0$ donc $f$ est strictement décroissante;
\item sur $\left ] \ln(3)\;;\;+\infty\strut\right [$, $f'(x)> 0$ donc $f$ est strictement croissante;
\item la fonction $f$ admet pour $x=\ln(3)$ un minimum.
\end{list}

		\item 
\begin{list}{\textbullet}{}
\item $\ln(3)\approx 1,1$ donc $[2~;~5] \subset\; ]\ln(3)\;;\;+\infty[$, et donc la fonction $f$ est strictement croissante sur $[2~;~5]$.
\item $f(2)=\ln\left (1+\e^{-2}\right ) + \dfrac{1}{4}\times 2 \approx 0,63$ donc $f(2)<1$.
\item $f(5)=\ln\left (1+\e^{-5}\right ) + \dfrac{1}{4}\times 5 \approx 6,25$ donc $f(5)>1$.	\item De plus, la fonction $f$ est dérivable sur $\R$ donc continue.
\end{list}
	
D'après le corollaire du théorème des valeurs intermédiaires, 	on peut déduire que l'équation $f(x)=1$ admet une solution unique sur $[2~;~5]$. On l'appelle $\alpha$.
	\end{enumerate}
	
\end{enumerate}

\medskip

\textbf{Partie B}

\medskip

On admettra que la fonction $f'$ est dérivable sur $\R$ et pour tout réel $x$,
$f''(x) = \dfrac{\e^x}{\left(\e^x + 1\right)^2}$.

On note $\Delta$ la tangente à la courbe $\mathcal{C}_f$ au point d'abscisse $0$.

Dans le graphique ci-dessous, on a représenté la courbe $\mathcal{C}_f$ la tangente $\Delta$ et le quadrilatère MNPQ tel que M et N sont les deux points de la courbe $\mathcal{C}_f$ d'abscisses respectives $\alpha$ et $- \alpha$, et Q et P sont les deux points de la droite $\Delta$ d'abscisses respectives $\alpha$ et $- \alpha$.


\psset{unit=1.2cm,arrowsize=2pt 3}
\begin{center}
\begin{pspicture}(-6,-1)(6,5)
\pspolygon[linewidth=1.25pt,linecolor=blue,fillstyle=solid,fillcolor=blue!20](3.922,-0.3)(3.922,1)(- 3.922,2.961)(-3.922,1.7)

\psaxes[linewidth=1.25pt,Dx=10,Dy=10](0,0)(-6,-1)(6,5)
\psaxes[linewidth=1.25pt,Dx=10,Dy=10]{->}(0,0)(1,1)
\psplot[plotpoints=2000,linewidth=1.25pt,linecolor=red]{-6}{6}{1 1 2.71828 x exp div add ln 0.25 x mul add}
\psplotTangent[linewidth=1.25pt,linecolor=blue]{0}{6}{1 1 2.71828 x exp div add ln 0.25 x mul add}

\psline[linestyle=dashed](-3.922,1.7)(-3.922,0)
\uput[d](0.5,0){$\vect{\imath}$}\uput[l](0,0.5){$\vect{\jmath}$}
\uput[d](3.922,-0.3){Q}\uput[u](3.922,1){M}\uput[ur](- 3.922,2.961){N}\uput[dl](-3.922,1.7){P}
\uput[dr](3.922,0){$\alpha$}\uput[d](-3.922,0){$- \alpha$}\uput[dl](0,0){O}\uput[d](-5.2,2){\blue$\Delta$}
\uput[d](-5.2,3.8){\red $\mathcal{C}_f$}
\end{pspicture}
\end{center}

\smallskip

\begin{enumerate}
\item 
	\begin{enumerate}
		\item 
\begin{list}{\textbullet}{Pour tout réel $x$:}
\item $\e^{x}>0$;
\item $\e^{x}>0$ donc $\e^{x}+1>0$ donc $\left(\e^x + 1\right)^2>0$.
\end{list}		

Donc $f''(x) = \dfrac{\e^x}{\left(\e^x + 1\right)^2}>0$.

		\item
Pour tout réel $x$, $f''(x)>0$ donc la fonction $f$ est convexe; sa courbe représentative est donc située au-dessus de toutes ses tangentes, donc au-dessus de $\Delta$.

De plus, comme la fonction $f$ est convexe, sa courbe représentative entre les points M et N est située en-dessous de la sécante (MN).

Donc  la portion de la courbe $\mathcal{C}_f$ sur l'intervalle $[-\alpha~;~\alpha]$, est inscrite dans le quadrilatère MNPQ.
	\end{enumerate}
	
\item 	
	\begin{enumerate}
		\item 
$\begin{gathered}f(-\alpha) = \ln\left (1+\e^{\alpha}\right ) -\dfrac{1}{4}\alpha
= \ln \left (\e^{\alpha} \left (\e^{-\alpha}+1\right )\right ) -\dfrac{1}{4}\alpha
= \ln\left (\e^{\alpha}\right ) + \ln \left (\e^{-\alpha}+1\right ) -\dfrac{1}{4}\alpha\\[7pt]
\phantom{f(-\alpha)}
= \alpha + \ln \left (\e^{-\alpha}+1\right ) -\dfrac{1}{4}\alpha
= \ln \left (\e^{-\alpha}+1\right )+ \dfrac{3}{4}\alpha\end{gathered}$
		
		\item 
On sait que $f(\alpha)=1$, donc $\ln\left (1+\e^{-\alpha}\right )+\dfrac{1}{4}\alpha = 1$ et donc $\ln\left (1+\e^{-\alpha}\right ) = 1-\dfrac{1}{4}\alpha$.

Donc comme $f(-\alpha) = \ln \left (\e^{-\alpha}+1\right )+ \dfrac{3}{4}\alpha$, on a:
$f(-\alpha) = 1-\dfrac{1}{4}\alpha + \dfrac{3}{4}\alpha = 1+\dfrac{1}{2}\alpha$.


Les points M et N ont respectivement pour coordonnées $\left (\alpha\;;\;f(\alpha)\right )$ et $\left (-\alpha\;;\;f(-\alpha)\right )$ donc le coefficient directeur de la droite (MN) est égal à:

\[\dfrac{y_{\text M} - y_{\text N}}{x_{\text M} - x_{\text N}}
= \dfrac{f(\alpha) - f(-\alpha)}{\alpha-(-\alpha)}
= \dfrac{1-\left (1+\frac{1}{2}\alpha\right )}{2\alpha}
= \dfrac{-\frac{1}{2}\alpha}{2\alpha}
= -\dfrac{1}{4}\]

La droite $\Delta$ est tangente en 0 à la courbe $\mathcal{C}_f$; donc son coefficient directeur est égal à:
$f'(0)=\dfrac{\e^{0}-3}{4\left (1+\e^{0}\right )}
= \dfrac{1-3}{4\left (1+1\right )}
= \dfrac{-2}{8}
= -\dfrac{1}{4}$

Les droites (MN) et $\Delta$ ont le même coefficient directeur, donc elles sont parallèles, donc $(\text{MN}) /\!/ (\text{QP})$.

Par construction, les droites (MQ) et (NP) sont parallèles.

On peut donc en déduire que le quadrilatère MNPQ est un parallélogramme.
	\end{enumerate}
\end{enumerate}
